\section{Conclusions}
We have shown that GCL is not restricted to only link perturbations or feature augmentation.
%we attempt to answer a question: are we really  restricted to the data augmentation for the GCL? 
By introducing a simple and efficient spectral feature augmentation layer %realized by the fast incomplete power iteration, %between GNN encoder and projection heads, we demonstrate that our feature augmentation strategy improves the performance of  GCL.  
%
we achieve significant performance gains. 
Our incomplete power iteration is very fast. % \ie, for low $k$, the cost of SFA is negligible (see the \nameref{sec:timing} section). % what happened beyond the method. 
%By subtracting a low-rank approximation of features form themselves, we obtain the augmented features. 
%
%
Our theoretical analysis has demonstrated that SFA rebalances the useful part of spectrum, and also augments the useful part of spectrum by implicitly injecting the noise into singular values of both data views. SFA leads to  a better alignment with a lower generalization bound. 
%
%We concluded that there are several  interesting properties for such augmentation strategy. Firstly, it  adds the noise to the spectrum of the feature matrix. Secondly, it promotes  the orthonormal bases  alignment between features in two views. Thirdly, we establish the relation between the spectral noise and the feature noise.
%
%In experiments, we demonstrate the competitive performance compared with other baselines. In addition, we also found that our method seems to alleviate the dimension collapse, as our smallest singular values are non-zero. Importantly, our method is compatible with existing data augmentation.
